%% =====================================================================
%% 这是 main.tex —— 南昌大学本科生毕业设计（论文）模板示例文档
%% 编译：xelatex -> bibtex -> xelatex -> xelatex
%% =====================================================================
\documentclass{ncubachelor}
%% 常用类选项（可叠加，方括号内逗号分隔）：
%%   twoside        双面打印：奇偶页眉在外侧分别显示章名（复旦思路），
%%                  并启用对称的内外侧边距与“空偶数页不显示页眉”
%%   blindreview    盲审版：封面/摘要中的姓名、学号、导师自动替换为 ***，
%%                  并清除 PDF 元数据中的作者信息
%%   fontset=...    字体方案：auto(默认,自动探测)/windows/mac/founder/adobe/fandol/none
%%   tocdots        目录显示点线（v1.1.5 起默认开启，tocdots=false 关闭）
%%   openright      每章从奇数页开始（配合 twoside 用于双面印刷）
%% 示例：\documentclass[twoside,blindreview]{ncubachelor}

%% ---------------------------------------------------------------------
%% 一、论文元信息（按实际填写）
%% ---------------------------------------------------------------------
\ncusetup{
  info = {
    secret-level     = {公开},
    title            = {南昌大学本科生毕业论文模板使用演示},
    title*           = {Demonstration of the Nanchang University\\Undergraduate Thesis Template},
    college          = {经管},
    department       = {工商},
    major-class      = {物理学 2023 班},
    major            = {物理学},
    author           = {张三},
    studentid        = {402230220001},
    supervisor       = {李四},
    supervisor-title = {教授},
    start-date       = {23},
    stop-date        = {27},
    start-stop-date  = {2023 年 9 月—2027 年 6 月}
  }
}

%% ---------------------------------------------------------------------
%% 二、模板细节配置（可选，默认值如下）
%% ---------------------------------------------------------------------
%% 盲审细节（配合 blindreview 类选项）：
%%   blind = {
%%     text = {***},                  %% 遮蔽占位文本
%%     hide-acknowledgement = true,   %% 同时省略致谢页
%%   }
%% 定理环境样式：
%%   theorem = {
%%     header-font = \songti\bfseries,  %% 头部字体
%%     within      = chapter,           %% 编号锚点：chapter / section / none
%%   }
%%   proof = { qed = $\blacksquare$ }  %% 证明结束符

%% ---------------------------------------------------------------------
%% 三、其他宏包与自定义设置（按需添加）
%% ---------------------------------------------------------------------
% \usepackage{siunitx}
% \usepackage{algorithm,algorithmicx,algpseudocode}
% \usepackage{listings}

%% ---------------------------------------------------------------------
%% 三、文档开始
%% ---------------------------------------------------------------------
\begin{document}

%% 封面（无页眉页脚）
\maketitle

%% 原创性申明与授权书（无页眉页脚）
\makedecaut

%% 前置部分：页码从中文摘要起用罗马数字 I, II, III …
\frontmatter

%% 中文摘要
\begin{abstract}
这是中文摘要。摘要应扼要说明研究目的、方法、结果和结论，
内容为小四号宋体、两端对齐、行距 1.35 倍、行首缩进 2 个汉字字符。

如果摘要较长，可以分段书写，也可以写到第二页。

本文得到了国家自然科学基金项目的资助。
\keywords{南昌大学；毕业论文；\LaTeX{} 模板}
\end{abstract}

%% 英文摘要
\begin{abstract*}
This is the English abstract. It should state briefly the purpose,
methods, results and conclusions of the research.

If the abstract is long enough, it may occupy more than one page.

This work was supported by the National Natural Science Foundation
of China.
\keywords{Nanchang University; bachelor thesis; \LaTeX{} template}
\end{abstract*}

%% 目录
\tableofcontents

%% ---------------------------------------------------------------------
%% 正文：页码从第一章起改用阿拉伯数字 1, 2, 3 …
%% ---------------------------------------------------------------------
\mainmatter

\chapter{引言}\label{chap:intro}

%% 若论文源于科研项目，可在正文第一页用无编号脚注标注项目来源
\blfootnote{本论文来源于科研项目“基于 \LaTeX{} 的本科毕业设计”，项目编号 1124。}

\section{研究背景与意义}\label{sec:background}

我们知道 lshort\cite{lshortcn} 是入门 \LaTeX{} 的必读文献。
可以花上几个小时阅读一下，能够避免被一些简单的问题耽误时间。
除此之外，还推荐刘海洋的《\LaTeX{} 入门》。

在社交媒体的强覆盖下，新闻信息的传播渠道也悄然发生了变化。
根据美国皮尤研究中心的调查结果\cite{pew}，社交媒体正在成为网络上
热点事件生成并发酵的源头，国内较为知名的新浪微博\footnote{本脚注
为小五号宋体，仅作排版示例。}即是典型代表。

\section{本文工作}

本文的主要工作如下：
\begin{enumerate}
  \item 依据《南昌大学本科生毕业设计（论文）书写式样》实现了
        封面、申明与授权书、中英文摘要、目录、正文、参考文献与致谢
        的自动排版；
  \item 提供了三线表、定理类环境、脚注等常用排版要素；
  \item 通过示例演示了模板的用法。
\end{enumerate}

%% ---------------------------------------------------------------------
\chapter{图表与公式}\label{chap:figtab}

\section{表格}

表格标题置于表的上方，五号宋体加粗居中，表序与表名之间空一个汉字宽度；
表格内容为五号宋体，行距 1.35。表~\ref{tab:demo} 给出了一个三线表示例。

\begin{table}[htb]
  \centering
  \caption{用不同技术得到的带隙温度系数、$E_{g0}$、$\alpha$ 和 $T_0$ 的值}
  \label{tab:demo}
  \begin{tabular}{llccccc}
    \thickhline
    样品类型 & 实验方法 & $\mathrm{d}E_g/\mathrm{d}T$ (eV/K) &
    $E_{g0}$(eV) & $\alpha$ (eV/K) & $T_0$(K) & 文献 \\
    \hline
    GaN/Al$_2$O$_3$ & 光致发光 & $-5.32\times 10^{-4}$ & 3.503 &
    $5.08\times 10^{-4}$ & $-996$ & 61 \\
    GaN/Al$_2$O$_3$ & 光致发光 & —                      & 3.489 &
    $7.32\times 10^{-4}$ & 700   & 59 \\
    GaN/Al$_2$O$_3$ & 光吸收   & $-4.5\times 10^{-4}$   & $\sim$3.471 &
    $-9.3\times 10^{-4}$ & 772   & 63 \\
    \thickhline
  \end{tabular}
\end{table}

\section{插图}

图标题置于图的下方，五号宋体加粗居中，图序与图名之间空一个汉字宽度。
图~\ref{fig:demo} 给出了示例。请注意，浮动体不一定会留在插入位置，
因此不要使用“下图”“下表”这类表述，应使用“如图 1-1 所示”。

\begin{figure}[htb]
  \centering
  \includegraphics[width=.55\textwidth]{example-image}
  \caption{热风速计原理}
  \label{fig:demo}
\end{figure}

\section{公式}

对 Fourier 变换采用以下约定：
\begin{equation}
  (\mathcal{F}f)(\xi) = \hat{f}(\xi)
    = \frac{1}{\sqrt{2\pi}}\int_{-\infty}^{\infty}
      \mathrm{e}^{-ix\xi}f(x)\,\mathrm{d}x.
  \label{eq:fourier}
\end{equation}
根据这一标准定义，有
\[
  \lVert\hat{f}\rVert_{L^2} = \lVert f\rVert_{L^2},
  \qquad
  |\hat{f}(\xi)| \le (2\pi)^{-1/2}\lVert f\rVert_{L^1}.
\]
式~\eqref{eq:fourier} 中的积分顺序可以互换。

\section{定理类环境}

\begin{definition}[排版系统]
排版系统是把文字、公式、图表按规范组织成文档的软件系统。
\end{definition}

\begin{theorem}[可编译性]
只要元信息填写完整，本模板即可通过 XeLaTeX 顺利编译。
\end{theorem}

\begin{proof}
即得易见平凡，仿照上例显然。留作习题答案略，读者自证不难。
反之亦然同理，推论自然成立。略去过程 QED，由上可知证毕。
\end{proof}

\begin{lemma}
模板还提供引理、推论、命题、公理、假设、性质、注解等环境，
编号格式统一为“章.序”。
\end{lemma}

\begin{example}
在 \LaTeX{} 中书写“例”环境：本段文字即“例”环境的正文内容。
\end{example}

\begin{remark}
定理与证明的头部字体、编号锚点、证明结束符均可通过
\verb|\ncusetup{theorem=...}| 与 \verb|\ncusetup{proof=...}| 调整。
\end{remark}

%% ---------------------------------------------------------------------
\unnumberedchapter{结论}

结论一般不参与章节编号，但要进目录并出现在页眉中，
因此这里使用无编号章命令 \verb|\unnumberedchapter{结论}|。

本章总结全文的主要工作与不足，并对后续研究进行展望。

%% ---------------------------------------------------------------------
\appendix
\chapter{附加材料}\label{chap:appendix}

附录与正文的章节标题用法没有区别，仅编号方式不同（附录 A、附录 B）。

%% ---------------------------------------------------------------------
%% 后置部分
%% ---------------------------------------------------------------------
\backmatter

%% 参考文献
\bibliography{references}

%% 致谢
\begin{acknowledgements}
感谢……
\signoff{张三}{2027 年 6 月}
\end{acknowledgements}

\end{document}
%% 结束
